\section{Stochastic Processes}

\answer{RS1}
Use four nonterminal prefixes: the empty prefix, H, HT and HTH. Write their remaining means as $e_0,e_1,e_2,e_3$. When H follows HTH, the longest useful suffix is H, not the empty string. First-step equations are
\[
 e_0=2+e_1,\quad e_1=2+e_2,\quad
 e_2=1+\tfrac12e_3+\tfrac12e_0,\quad
 e_3=1+\tfrac12e_1.
\]
The first two equations give $e_1=e_0-2,e_2=e_0-4$, and the last gives $e_3=e_0/2$. Substitution yields $e_0-4=1+3e_0/4$, hence $\boxed{e_0=20}$. Each block of four tosses equals HTHT with probability $1/16$, ensuring finite mean.

\answer{RS2}
The stationary geometric ratio solves $\tfrac12r^2-r+\tfrac14=0$. Its root in $(0,1)$ is $r=1-1/\sqrt2$, so
\[
 \boxed{\Prob_\pi(X\ge2)=\left(1-\frac1{\sqrt2}\right)^2},\qquad
 \boxed{\E_0\tau_0^+=\frac1{1-r}=\sqrt2}.
\]
Clearing or attempting service on an empty queue counts as a positive-time return, so the mean can be less than 2.

\answer{RS3}
A two-step switch occurs exactly when one of the two steps switches. Thus $2\alpha(1-\alpha)=3/8$, giving
\[
 \boxed{\alpha=1/4\quad\text{or}\quad\alpha=3/4}.
\]
Both parameters give the same two-step transition matrix. With the same initial law, all joint distributions at even times are consequently identical, so such data cannot distinguish them. One-step transition information can: one model favors persistence and the other favors alternation.

\answer{RS4}
Weights 3 and 2 for pending A and B tasks decrease in conditional expectation by 1 per processed task, proving a finite expected total. The means satisfy $a=1+b,b=1+a/3$, giving $a=3,b=2$. If second moments are $u,v$, independence of the two B subtrees gives
\[
 u=1+2b+v+b^2=9+v,\qquad
 v=1+\tfrac23a+\tfrac13u=3+\tfrac13u.
\]
These recursions may first be applied to trees truncated by generation. Their means are bounded by 3 and 2, and the second-moment recursion is bounded by 18 and 9. Monotone convergence proves finite second moments. Solving then yields
\[
 \boxed{\E T=3,\quad\E T^2=18,\quad\Var T=9}.
\]

\answer{RH1}
In R11, $\rho=1/2$ gives $r=2-\sqrt3$, hence
\[
 \boxed{\E T=2+2\sqrt3}.
\]
Whether step $t$ must be performed is determined by the past and is independent of its reset decision. Therefore the total reset count $K$ satisfies $\E K=(1/2)\E T=\boxed{1+\sqrt3}$. The pathwise cost is $T+2K$, giving
\[
 \boxed{\E C=4+4\sqrt3}.
\]

\answer{RH2}
The drawdown gap has a self-loop at zero, so a two-absorbing-boundary time formula does not apply. In R16, $d=3,D=12$, giving
\[
 \boxed{\E\tau=12,\qquad\Var\tau=92},\qquad
 \boxed{\Prob(M_\tau\ge4)=\left(\frac34\right)^4=\frac{81}{256}}.
\]
Each increase in the maximum requires winning a fresh boundary race with distances 1 and 3.

\answer{RH3}
The denominator in R13 is $D=6(1/4)+1/2=2$, so
\[
 \boxed{h_3^+=\frac{3/4+1/2}{2}=\frac58},\qquad
 \boxed{h_3^-=\frac{3/4}{2}=\frac38}.
\]
With an unobserved equally likely initial direction, average the two conditional laws to obtain success probability $\boxed{1/2}$. This averaging does not automatically make the position process memoryless.

\answer{RH4}
The probabilities that the first 0, 1, 2 or 3 draws are distinct are $1,1,2/3,2/9$. Thus
\[
 \E N=1+1+\frac23+\frac29=\frac{26}{9}.
\]
The mean cost per draw is $7/3$. The past-measurable continuation-event argument of R18 gives
\[
 \boxed{\E C=\frac73\frac{26}{9}=\frac{182}{27}}.
\]
There are exactly $N-1$ distinct types before the first repeat, and the last draw adds none. Their mean count is therefore $\boxed{17/9}$.

\answer{RC1}
The up and down payoffs are 0 and $1/2$. Exact replication solves
\[
 2\Delta+B=0,\qquad\tfrac12\Delta+B=\tfrac12,
\]
giving $\boxed{\Delta=-1/3,B=2/3,V=1/3}$. If $\Delta\ge0$, covering the down state requires $\tfrac12\Delta+B\ge1/2$, hence initial cost $\Delta+B\ge1/2+\Delta/2\ge1/2$. Cash $1/2$ covers both outcomes, so $\boxed{V^{\rm super}=1/2}$.

\answer{RC2}
After an initial H, the next H probability is $2/3$, so the expected gain from continuing is $1/3$. After an initial T, it is $-1/3$, so stop. Toss initially, continue after H and stop after T. The initial value is
\[
 \boxed{V=\tfrac12(1+\tfrac13)+\tfrac12(-1)=\frac16}.
\]
Stopping immediately gives zero and is inferior. The second-stage comparison includes all feasible actions.

\answer{RC3}
Replacement pays net $7/8-X$, so replace when $X<7/16$, with either action at equality. Put $a=7/16$. Integration yields
\[
 \boxed{V=\int_0^a(7/8-x)\dd x+\int_a^1x\dd x
 =\frac12+\frac{49}{256}=\frac{177}{256}}.
\]
The constant continuation value of an independent redraw does not apply; here the replacement value depends on the observed $X$.

\answer{RC4}
For all-pass testing, the indices are $c/(1-p)=(2,8,4/3)$. The order is $(3,1,2)$, with cost
\[
 \boxed{1+\tfrac14\cdot1+\tfrac14\cdot\tfrac12\cdot2=\frac32}.
\]
For at-least-one-pass testing, the indices are $c/p=(2,8/3,4)$. The order is $(1,2,3)$, with cost
\[
 \boxed{1+\tfrac12\cdot2+\tfrac12\cdot\tfrac14\cdot1=\frac{17}{8}}.
\]
Each later cost is weighted by the probability that earlier tests have not yet settled the decision.

\answer{RG1}
The variables $1,W_1,W_1^2-1$ are orthogonal in square-integrable space; the last two have variances 1 and 2. R25 gives $\E I=1$, $\Cov(I,W_1)=1/2$ and $\Var I=2e-5$. Differentiating normal exponential moments also gives
\[
 \E[e^{W_t-t/2}(W_1^2-1)]=t^2,\qquad
 \Cov(I,W_1^2-1)=\int_0^1t^2\dd t=\frac13.
\]
Thus the best predictor and residual error are
\[
 \boxed{\widehat I=1+\frac12W_1+\frac16(W_1^2-1)}.
\]
\[
 \boxed{\E(I-\widehat I)^2=2e-5-\frac14-\frac1{18}=2e-\frac{191}{36}}.
\]
The quadratic term reduces error by another $1/18$, although a finite-degree polynomial still differs from the full conditional mean.

\answer{RG2}
Substitute $T=2,b=0,a=1$ into R27. At time 1, the mean is $6/8=3/4$ and the variance is $1-1/2-3/8=1/8$, so
\[
 \boxed{W_1\mid(W_2=0,A=1)\sim N(3/4,1/8)}.
\]
The two other times have conditional covariance
\[
 \boxed{\frac12-\frac{(1/2)(3/2)}2
 -\frac{3(1/2)(3/2)(3/2)(1/2)}8=-\frac{11}{128}}.
\]
The area constraint makes their deviations partly compensate each other.

\answer{RG3}
The conditional reflection formula with $T=2,b=0$ gives maximum tail $\Prob(M\ge a)=e^{-a^2}$ for $a\ge0$. Therefore
\[
 \boxed{\Prob(M\ge1\mid W_2=0)=e^{-1}},\qquad
 \boxed{\E(M\mid W_2=0)=\int_0^\infty e^{-a^2}\dd a=\frac{\sqrt\pi}{2}}.
\]

\answer{RG4}
In R28, $k=\sqrt{2\lambda}=1$, giving
\[
 \boxed{u(1)=\frac{\sinh1}{\sinh2}=\frac1{2\cosh1}},\qquad
 \boxed{\E(\tau\mid\text{success})=2\coth2-\coth1=\tanh1}.
\]
The simplification uses $\sinh2=2\sinh1\cosh1$ and $\cosh2=\cosh^21+\sinh^21$. Without failure, success has probability $1/2$ and conditional mean duration $(4-1)/3=1$. Failure lowers the success probability and selects faster paths among those that succeed.

\par\Needspace{8\baselineskip}
\subsection*{Section Exercises}
\answer{OR01}
Partition the tosses into disjoint blocks of length $n$. Each block consists entirely of heads with probability $2^{-n}$, independently of the other blocks. Thus $T_n$ is bounded by $n$ times a geometric random variable and has finite mean.

Let $e_k$ be the expected number of additional tosses when the current run of heads has length $k$, for $0\le k\le n$. The terminal condition is $e_n=0$. A tail resets the run to zero, so
\[
 e_k=1+\tfrac12e_{k+1}+\tfrac12e_0,
 \qquad 0\le k<n.
\]
Substitution backwards from $e_n=0$ gives
\[
 e_0=\left(1+\frac{e_0}{2}\right)
       \sum_{j=0}^{n-1}2^{-j}
     =\left(1+\frac{e_0}{2}\right)(2-2^{1-n}).
\]
Rearranging yields $2^{-n}e_0=2-2^{1-n}$ and therefore
\[
 \boxed{\E T_n=2^{n+1}-2}.
\]
A run requires retaining the current progress: a tail destroys all consecutive heads accumulated so far.

\answer{OR07}
First suppose $0<p<1$ and write $q=1-p$. Let $X_t$ count the umbrellas at the departure location for trip $t$. The state space is $\{0,1,\ldots,n\}$. The next trip departs from the other location, so
\[
 P(0,n)=1,\qquad
 P(i,n-i)=q,\quad P(i,n-i+1)=p\quad(i\ge1).
\]
This finite chain is irreducible. For $i<n$, a trip without an umbrella followed by a trip carrying one reaches $i+1$. For $i\ge2$, the reverse sequence reaches $i-1$; from 1, the chain can also reach 0 through $n$. All these paths have positive probability.

Consider
\[
 \pi_0=\frac q{n+q},\qquad
 \pi_i=\frac1{n+q}\quad(1\le i\le n).
\]
The inflow to 0 is $q\pi_n=\pi_0$. For $1\le j<n$, the inflow is $q\pi_{n-j}+p\pi_{n-j+1}=\pi_j$; the inflow to $n$ is $\pi_0+p\pi_1=\pi_n$. The probabilities sum to 1, so this is the unique stationary distribution.

Getting wet requires state 0 and rain on the current trip. The long-run reward theorem for a finite irreducible chain therefore gives the almost-sure limiting frequency
\[
 \boxed{p\pi_0=\frac{p(1-p)}{n+1-p}}.
\]
The limit does not depend on $k$. If $p=0$, no trip is wet. If $p=1$, only the first trip can be wet: every subsequent departure location has an umbrella. The limiting frequency is again 0, as the same formula predicts.

\answer{OR11}
Write $q=1-p$ and $u=u_1$. Since a downward jump has size 1, a walk from $n$ to 0 must first visit $n-1,\ldots,1$ in order. Applying the strong Markov property and translation invariance at these descending hitting times gives
\[
 u_n=u^n.
\]
There may be arbitrarily many rightward jumps between successive descending records; the leftward jumps need not be consecutive. Conditioning on the first step from 1 gives
\[
 u=p+qu_3=p+qu^3,
\]
and hence
\[
 (u-1)(qu^2+qu-p)=0.
\]
The positive root of the quadratic factor is
\[
 r=\frac{\sqrt{1+4p/q}-1}{2}.
\]
If $p\ge2/3$, then $r\ge1$ and the other quadratic root is negative. Thus the only possible value in $[0,1]$ is 1.

If $p<2/3$, the mean increment is $2-3p>0$, so the strong law implies that the walk tends to positive infinity almost surely. Consequently $u\ne1$. Indeed, if every integer's left neighbor were reached almost surely, repeated use of the strong Markov property would imply that every negative integer is visited almost surely, contradicting convergence to positive infinity. We must therefore choose $u=r<1$. Thus
\[
 \boxed{u_n=
 \begin{cases}
 \left(\dfrac{\sqrt{1+4p/(1-p)}-1}{2}\right)^n,
       &0<p<2/3,\\[4pt]
 1,&2/3\le p<1.
 \end{cases}}
\]
A first-step equation may have multiple roots in the probability range. Drift, or the minimal nonnegative solution property, identifies the correct root.

\answer{OR16}
From every unabsorbed state, $a+b$ consecutive rightward steps force absorption. Such a block has probability $2^{-(a+b)}$. Hence $\tau$ has a geometric tail and all its moments are finite.

The symmetric independent increments make the following processes martingales:
\[
 S_n,\qquad S_n^2-n,\qquad S_n^3-3nS_n,
\]
\[
 S_n^4-6nS_n^2+3n^2+2n.
\]
These identities follow by expanding $(S_n\pm1)^j$ and averaging. Apply optional stopping first at the bounded time $\tau\wedge N$, then let $N\to\infty$. The stopped position is bounded in absolute value by $\max(a,b)$, and $\tau$ has finite moments, so dominated convergence justifies the limit.

Set $r=\Prob(S_\tau=b)$. The first two martingales give
\[
 br-a(1-r)=0,\qquad
 \E\tau=\E S_\tau^2=b^2r+a^2(1-r).
\]
Therefore
\[
 \boxed{r=\frac a{a+b}},\qquad \boxed{\E\tau=ab}.
\]
To obtain the second moment, write $A=\E[\tau\ind_{\{S_\tau=b\}}]$ and $B=\E[\tau\ind_{\{S_\tau=-a\}}]$. The cubic martingale yields
\[
 A+B=ab,\qquad
 bA-aB=\frac{\E S_\tau^3}{3}=\frac{ab(b-a)}3.
\]
Solving these two linear equations gives
\[
 A=\frac{ab(2a+b)}{3(a+b)},\qquad
 B=\frac{ab(a+2b)}{3(a+b)}.
\]
Consequently,
\[
 \E(\tau S_\tau^2)=b^2A+a^2B
       =\frac{ab(a^2+ab+b^2)}3,
\]
\[
 \E S_\tau^4=ab(a^2-ab+b^2).
\]
Substitute these values into the stopped quartic identity
\[
 0=\E S_\tau^4-6\E(\tau S_\tau^2)
       +3\E\tau^2+2\E\tau
\]
to obtain
\[
 \boxed{\E\tau^2=\frac{ab(a^2+3ab+b^2-2)}3}.
\]

\answer{OR20}
Let $V(a,b)$ be the maximum expected additional score when $a$ red and $b$ black cards remain. Stopping immediately adds 0, and the boundary values are
\[
 V(a,0)=a,\qquad V(0,b)=0.
\]
For $a,b>0$, the value of drawing one more card is
\[
 C(a,b)=\frac{a[1+V(a-1,b)]+b[-1+V(a,b-1)]}{a+b}.
\]
The Bellman equation is therefore
\[
 \boxed{V(a,b)=\max\{0,C(a,b)\}}.
\]
Evaluate states in increasing order of $a+b$. Backward induction proves optimality: once the successors use optimal policies, continuing is worth exactly $C(a,b)$, and comparison with the stopping value 0 gives the best current decision. The deck shrinks after every draw, so there is no infinite-horizon or integrability issue.

Starting with 26 cards of each color, the current cumulative score is $s=b-a$ and the number remaining is $N=a+b$. One optimal rule is as follows. Stop if $N=0$. Continue if $N>0$ and $s\le0$. For positive scores, continue precisely under the conditions in the table; otherwise stop:
\[
 \begin{array}{c|ccccc}
 s&1&2&3&4&5\\\hline
 \text{Continue if}&N\ge5&N\ge10&N\ge19&N\ge30&N\ge43
 \end{array}
\]
Stop whenever $s\ge6$. Only reachable states are included, so $N$ and $s$ have the same parity. At $s=1,N=3$, continuing and stopping both have value 0; the table chooses to stop.

Exact rational evaluation of the recurrence gives
\[
 \boxed{V(26,26)=\frac{41984711742427}{15997372030584}
       \approx2.624475549}.
\]
\par\smallskip\noindent\begin{minipage}{\linewidth}
The recurrence also determines the decision comparison at every state. The following code computes the exact initial value without rounding:
\begin{verbatim}
from fractions import Fraction as F
V = {}
for a in range(27):
    for b in range(27):
        if a == 0: V[a,b] = F(0)
        elif b == 0: V[a,b] = F(a)
        else:
            c = (a*(1+V[a-1,b])
                 + b*(-1+V[a,b-1]))/(a+b)
            V[a,b] = max(F(0), c)
print(V[26,26])
\end{verbatim}
\end{minipage}\par\smallskip
The advantage comes from information about the remaining deck. The successive draws are not independent gambles with a fixed success probability of $1/2$.

\answer{OR22}
Let $V$ be the optimal value before the initial free roll. At an observed face $j$, stopping pays $j$, while restarting has value $V-1$. The candidate Bellman equation is
\[
 V=\frac16\sum_{j=1}^6\max\{j,V-1\}.
\]
The value $V=4$ satisfies it, since
\[
 \frac{3+3+3+4+5+6}{6}=4.
\]
To verify an upper bound for all strategies, let $V_m$ be the value with at most $m$ rolls. We have $V_1=7/2\le4$ and
\[
 V_{m+1}=\frac16\sum_{j=1}^6\max\{j,V_m-1\}.
\]
The right side is increasing in $V_m$, so induction gives $V_m\le4$. Consider a strategy stopping at roll $N$ with $\E N<\infty$, and force it to stop by roll $m$. The resulting payoff converges to the original payoff. Faces are bounded and fees are controlled by $N$, so dominated convergence implies that its expected payoff is at most 4. If $\E N=\infty$, the final face is at most 6 and the expected net payoff is negative infinity, which cannot improve the value.

Thus reroll on 1 or 2, accept 4, 5, or 6, and take either action on 3. For example, stop at the first face at least 3. The accepted face is uniform on 3, 4, 5, 6, with mean $9/2$. The number of rolls is geometric with success probability $2/3$, so expected reroll fees are $3/2-1=1/2$. Hence
\[
 \boxed{\text{Maximum expected payoff}=9/2-1/2=4}.
\]
This policy terminates almost surely and attains the upper bound, proving optimality.

\answer{OR25}
Integration by parts for $tW_t$ gives
\[
 TW_T=\int_0^T W_t\,dt+\int_0^T t\,dW_t.
\]
Consequently,
\[
 X=\int_0^T(T-t)\,dW_t,\qquad
 Y=\int_0^T t\,dW_t.
\]
Both integrands are deterministic and square-integrable, so both integrals have mean 0. The It\^o isometry and its cross form give
\[
 \Var X=\int_0^T(T-t)^2dt=\frac{T^3}{3},\qquad
 \Var Y=\int_0^Tt^2dt=\frac{T^3}{3},
\]
\[
 \Cov(X,Y)=\int_0^Tt(T-t)dt=\frac{T^3}{6}.
\]
The cross form also follows by applying the isometry to the sum of the two integrands and expanding the square. Since $T>0$, both variances are nonzero. Therefore
\[
 \boxed{\Corr(X,Y)=\frac{T^3/6}{T^3/3}=\frac12}.
\]

\answer{OR26}
Define
\[
 \tau=\inf\{t\ge0:B_t^{(1)}=-1\}.
\]
The reflection principle gives $\Prob(\tau\le t)=2\Phi(-1/\sqrt t)$, where $\Phi$ is the standard normal distribution function. Hence $\tau<\infty$ almost surely, with density
\[
 f_\tau(t)=\frac1{\sqrt{2\pi t^3}}e^{-1/(2t)},\qquad t>0.
\]
The time $\tau$ depends only on the first coordinate and is independent of the entire second-coordinate process. Conditional on $\tau=t$, therefore, $B_\tau^{(2)}\sim N(0,t)$. Set $Z=B_\tau^{(2)}$. Mixing the normal densities yields
\[
 f_Z(z)=\int_0^\infty\frac1{2\pi t^2}
            e^{-(1+z^2)/(2t)}dt
       =\frac1{\pi(1+z^2)}.
\]
The integrand is nonnegative, so Tonelli's theorem justifies the mixture calculation; substitution $u=1/t$ evaluates the last integral. Thus $Z$ has the standard Cauchy distribution, and
\[
 \boxed{\Prob(1+Z>0)
   =\int_{-1}^\infty\frac{dz}{\pi(1+z^2)}
   =\frac34}.
\]
The independence is between the stopping time and the entire second-coordinate process. Evaluation at this random time produces the mixture above, not a normal distribution with fixed variance.

